%% kapa: exam=1 round=27 year=2016 period=1 form=A subject=economics subject_ko=경제학원론 seq=4 number=44
%% source: Q-Net 감정평가사 기출문제 (cst00309) artlSeq=5169402 file=extracted/1008/5169402/제27회 감정평가사 1차 문제/제27회 감정평가사 1교시문제(A형).pdf
%% conversion: kordoc 4.24.51 -> tools/gen1.py (auto); review=proofread-sonnet; flags=none
\begin{question}{44}
현재 소비자 甲은 주어진 소득 3,000원을 모두 사용하여 가격이 60원인 X재 20단위와 가격이 100원인 Y재 18단위를 소비하려고 한다. 이 때 X재와 Y재의 한계효용이 각각 20으로 동일하다면 효용극대화를 위한 甲의 선택으로 옳은 것은? (단, 소비자 甲의 X재와 Y재에 대한 무차별곡선은 우하향하고 원점에 대하여 볼록하다.)
\begin{choices}
  \item 현재 계획하고 있는 소비조합을 선택한다.
  \item X재 18단위와 Y재 18단위를 소비한다.
  \item X재 20단위와 Y재 20단위를 소비한다.
  \item X재의 소비량은 감소시키고 Y재의 소비량은 증가시켜야 한다.
  \item X재의 소비량은 증가시키고 Y재의 소비량은 감소시켜야 한다.
\end{choices}
\end{question}
\answer{5}
